\section{The ordered summation cover and its specialization}

The relevant cover must be named before giving its monodromy. Let $E/K$
be an elliptic curve, let $\pi:E\to\mathbb P^1$ be the quotient by
inversion, put $n=m-1$ with $m\geq3$, and take independent ordered
transcendentals $t_1,\ldots,t_n$. The corpus's theorem concerns
\[
 S_m(t_1,\ldots,t_n,T)\quad\text{over }L=K(t_1,\ldots,t_n),
\]
not an unspecified symmetric parameter quotient or every specialization
\cite{monodromySource}.

\begin{proposition}[Existing signed-sum structure statement]
Assume the quotient map $\pi$ is separable of degree two with deck
involution $[-1]\ne1$, and the signed sums are distinct at the generic
point up to simultaneous inversion. The root field and splitting field
are the fixed field of diagonal inversion in $K(E^n)$, and their group
over $L$ is
\[
 (\Z/2\Z)^n/\langle(1,\ldots,1)\rangle
       \simeq(\Z/2\Z)^{m-2},
\]
acting regularly on the $2^{m-2}$ roots.
\end{proposition}

The existing derivation obtains the product deck group from the
independent quotient maps. Its roots are
$x(\epsilon_1P_1+\cdots+\epsilon_nP_n)$. Diagonal inversion fixes all
roots. A further generic coincidence would force a nonempty proper
partial signed sum into $E[2]$; the preimage of that finite group scheme
under the corresponding surjective sum map is proper. This gives the
stabilizer exactly, and abelianness makes the single-root field the
splitting field. The corpus's longer manuscript also argues that the
constant field is unchanged, hence arithmetic and geometric monodromy
agree. A specialist review of the quotient and characteristic-two
scheme argument remains on its publication checklist; this synthesis
adds no independent proof review.

Over a finite field, exclude branch points, infinite signed sums, degree
drops and coincident roots. Frobenius on each degree-two fibre is a sign
bit. In the regular quotient action its image has order at most two.
Consequently a good specialization is either completely split or a
product of equal-degree irreducible quadratics. If all inputs are
abscissae of rational factor-base points, their sign bits are zero and
it is completely split. This last assertion also follows directly:
compute the signed sums of those rational points.

\paragraph{The exclusion this supplies.}
A factorization-type statistic is constant on this good rational
factor-base locus. Such a statistic cannot distinguish relation-rate
anomalies there. It does not follow that every statistic or algebraic
representation is constant, or that finding the other summands for a
fixed target is cheap. Complete splitting solves for the last coordinate
when the others are given; relation search asks a different question.

The older finding \art{KN-FIND-a1f3c2} had declared a broader open problem
closed. The later \art{KN-FIND-6edecd} explicitly declines that closure:
its separate toy corroboration addresses the full-symmetric-monodromy
premise, while the relation-rate and imprimitivity questions are not
closed by that empirical study. We therefore report the ordered-cover
structure and its stated exclusions, not the older universal closure
language. Permuting or symmetrizing the input coordinates changes the
base and can enlarge the monodromy; the elementary abelian description
must not be transferred to that different cover without a proof.
